Local PageRank methods connect sparse graph diffusion to numerical
optimization. The push algorithms of
\citet{andersen2006local,andersen2007using} charge neighborhood scans to
positive residual mass. The variational formulation of
\citet{fountoulakis2019variational} gives a weighted $\ell_1$ objective,
monotone local ISTA updates, and a bound on the optimum's support volume.
The regularization-path results of \citet{ha2021statistical} provide
a related foundation for continuation.

\paragraph{The closest RPPR algorithms.}
\citet{martinezrubio2023accelerated} introduced Conjugate Directions for
PageRank (CDPR), an exact method, and Accelerated Sparse PageRank (ASPR),
an approximate method, on expanding supports.
\citet{wei2026simple} replaced the restricted numerical solve by a
nearly-linear SDD solve, obtaining a bound proportional to the number
of support discoveries. Their output also supplies the positive-residual
certificate of Andersen, Chung, and Lang (ACL).
For the comparison in \cref{tab:lineage}, write $\mathcal S_\rho^*:=\supp(\vx_\rho^*)$ for the RPPR optimum's
support and let $k_*=|\mathcal S_\rho^*|$,
$V_*=\vol(\mathcal S_\rho^*)$, and
$\widetilde V_*=\nnz(\mQ_{\mathcal S_\rho^*,\mathcal S_\rho^*})$.
The first volume uses original graph degrees; the last quantity counts
internal nonzeros. \Cref{app:related-work} gives the exact normalization
maps, theorem pointers, and output qualifications.

\begin{table}[t]
\centering
\small
\caption{RPPR work in the canonical graph normalization. Additive
approximation refers to the RPPR objective gap. Soft bounds suppress the
logarithmic factors specified in the access model; probability qualifications
are shown explicitly. The comparison retains original-degree volume $V_*$
and internal nonzeros $\widetilde V_*$.}
\label{tab:lineage}
\begin{tabular}{@{}
 >{\raggedright\arraybackslash}p{0.35\textwidth}
 >{\raggedright\arraybackslash}p{0.23\textwidth}
 >{\raggedright\arraybackslash}p{\dimexpr0.42\textwidth-4\tabcolsep\relax}@{}}
\toprule
Method & Output & Work \\
\midrule
Local ISTA \citep{fountoulakis2019variational}
 & Additive approximation & $\softO(V_*/\alpha)$ \\
CDPR \citep{martinezrubio2023accelerated}
 & Exact optimum & $O(k_*^3+k_*V_*)$ \\
ASPR \citep{martinezrubio2023accelerated}
 & Additive approximation & $\softO(k_*\widetilde V_*/\sqrt\alpha+k_*V_*)$ \\
SDD active set \citep{wei2026simple}
 & Additive approximation and ACL certificate & $\softO(k_*V_*)$ with high probability \\
Threshold batching (this paper)
 & Additive approximation, safe support & $\softO(V_*\min\{k_*,\alpha^{-1/2}\})$ expected \\
Deterministic continuation (this paper)
 & Additive approximation, safe subsolution & $\softO(1/(\rho\sqrt\alpha))$ \\
\bottomrule
\end{tabular}
\end{table}

Wei and Yang also discuss a deterministic SDD substitution with
additional factors; the specific soft bound, rather than the existence
of a deterministic active-set method, is the distinction here.
The support-adaptive minimum also records that the existing
$\widetilde O(\rho^{-2})$ bound can be smaller in some regimes.

\paragraph{Classical acceleration and evolving sets.}
The local iterative methods of
\citet{zhou2024iterative} and the accelerated evolving
set framework of \citet{huang2025accelerated} develop the link between
numerical acceleration and changing local supports. Their guarantees and
accuracy dependence differ from the graph-uniform RPPR bound studied here.
\citet{fountoulakis2026complexity} give conditional guarantees for
classical accelerated iterations together with a support-locality
counterexample. Our result concerns the two algorithms specified in this
paper and preserves that distinction.

\Needspace{7\baselineskip}
\paragraph{Other approximation and access models.}
Some PPR estimators use the same output norm as our PPR corollary
\citep{wei2024absolute}; teleportation dependence, preprocessing, and access
assumptions also matter. Recent scalar PageRank estimation
\citep{thorup2026instance} and $\ell_1$-regularized resistance
\citep{li2026resistance} have different outputs or computational models.
\Cref{app:related-work} gives detailed comparisons with these and other
local diffusion, estimation, and obstacle methods.

\paragraph{Organization.}
\Cref{sec:problem-formulation} fixes the problem, access model, and main
results. \Cref{sec:obstacle-active-sets} develops the common geometry and
local certificates. \Cref{sec:deterministic-method} presents deterministic
continuation, its two-energy analysis, and its sparse implementation.
\Cref{sec:randomized-method} gives threshold batching and proves its depth
and total-work bounds. \Cref{sec:two-stage-ppr} develops the PPR consequences
and discusses further questions. Detailed literature comparisons, bounded
arithmetic, and seed extensions appear in the appendices.
